% Part: methods
% Chapter: proofs
% Section: example-1

\documentclass[../../../include/open-logic-section]{subfiles}

\begin{document}

\olfileid{mth}{prf}{ex1}

\olsection{একটি উদাহরণ}

প্রথম উদাহরণটি সেটের সংযোগ ও ছেদ নিয়ে একটি সরল
তথ্য। এতে সংজ্ঞা খুলে লেখা, সংযোজন ও সার্বিক
দাবি প্রমাণ এবং ক্ষেত্রবিচারে প্রমাণ দেখব।

\begin{prop}
যেকোনো সেট $A$, $B$ ও $C$-র জন্য
$A \cup (B \cap C) = (A \cup B)
\cap (A \cup C)$।
\end{prop}

চলো প্রমাণ করি!{}

\begin{proof}
যেকোনো সেট $A$, $B$ ও $C$-র জন্য
$A \cup (B \cap C) = (A \cup B) \cap (A \cup C)$
দেখাতে চাই।
\begin{quote}
প্রথমে প্রস্তাবের ``$=$''-র সংজ্ঞা খুলব। মনে করো,
দুটি সেট অভিন্ন প্রমাণ করতে দেখাতে হয় তাদের
!!{element}s একই। অর্থাৎ $A \cup (B \cap C)$-র
সব !!{element}s $(A \cup B) \cap (A \cup C)$-রও
!!{element}s, এবং উল্টোটিও। ``উল্টোটিও'' মানে
$(A \cup B) \cap (A \cup C)$-র প্রত্যেক
!!{element} $A \cup (B \cap C)$-রও
!!a{element}। তাই সংজ্ঞা খুলে বুঝলাম একটি
সংযোজন প্রমাণ করতে হবে। লিখে নিই:
\end{quote}
সংজ্ঞা অনুযায়ী, $A \cup (B \cap C) = (A \cup B)
\cap (A \cup C)$ তখনই, যখন $A \cup (B \cap C)$-র
প্রত্যেক !!{element} $(A \cup B) \cap (A \cup C)$-রও
!!a{element}, এবং $(A \cup B) \cap (A \cup C)$-র
প্রত্যেক !!{element} $A \cup (B \cap C)$-রও
!!a{element}।
\begin{quote}
এটি সংযোজন বলে দুটি অংশ আলাদা করে প্রমাণ
করতে হবে। প্রথমে দেখাই, $A \cup (B \cap C)$-র
প্রত্যেক !!{element} $(A \cup B) \cap (A \cup C)$-রও
!!a{element}।

এটি সার্বিক দাবি। তাই $A \cup (B \cap C)$-র
যেকোনো !!{element} নিয়ে দেখাব, সেটি
$(A \cup B) \cap (A \cup C)$-রও !!a{element}।
যেকোনো !!{element}টির নাম ধরি~$z$। প্রমাণ
এগোবে এভাবে:
\end{quote}
প্রথমে প্রমাণ করি, $A \cup (B \cap C)$-র প্রত্যেক
!!{element} $(A \cup B) \cap (A \cup C)$-রও
!!a{element}। $z \in A \cup (B \cap C)$ ধরি।
দেখাতে হবে $z \in (A \cup B) \cap (A \cup C)$।
\begin{quote}
এবার $\cup$ ও~$\cap$-র সংজ্ঞা খুলব। যেমন,
$\cup$-র সংজ্ঞা $A \cup B = \Setabs{z}{z \in A
  \text{ অথবা } z \in B}$। ``$A \cup (B \cap C)$''-তে
সংজ্ঞা প্রয়োগ করলে সংজ্ঞার ``$B$''-র ভূমিকায়
``$B \cap C$'' বসে। তাই $A \cup (B \cap C)
= \Setabs{z}{z \in A \text{ অথবা }
  z \in B \cap C}$। ফলে $z \in A \cup (B \cap C)$
পূর্বধারণাটির অর্থ $z \in \Setabs{z}{z \in A
\text{ অথবা } z \in B \cap
  C}$। আর $z \in \Setabs{z}{\dots z\dots}$
তখনই, যখন \dots $z$ \dots; এখানে অর্থাৎ
$z \in A$ অথবা $z \in B \cap C$।
\end{quote}
$\cup$-র সংজ্ঞায়, হয় $z \in A$ নয়তো
$z \in B \cap C$।
\begin{quote}
এটি বিযোজন বলে ক্ষেত্রবিচারে প্রমাণ কাজে লাগবে।
দুটি ক্ষেত্র নিয়ে প্রতিটিতে কাঙ্ক্ষিত সিদ্ধান্ত
``$z \in (A \cup B) \cap (A \cup C)$'' দেখাব।
\end{quote}
ক্ষেত্র 1: ধরো $z \in A$।
\begin{quote}
পূর্বধারণা থেকে আর বেশি কিছু পাওয়ার নেই। তাই
কাঙ্ক্ষিত সিদ্ধান্তটি দেখি। দেখাতে চাই
$z \in (A \cup B) \cap (A \cup C)$। $\cap$-র
সংজ্ঞা অনুযায়ী $z \in (A \cup B) \cap (A \cup C)$
দেখাতে হলে প্রমাণ করতে হবে, $z$
$(A \cup B)$ এবং $(A \cup C)$—দুটিতেই আছে।
কিন্তু $z \in A \cup B$ তখনই, যখন $z \in A$
অথবা $z \in B$; আর ক্ষেত্র~1-এর পূর্বধারণাতেই
$z \in A$ আছে। একই যুক্তিতে $B$-র জায়গায়
$C$ বসালে পাই $z \in A \cup C$। এই ভাবনাটি
উল্টো দিক থেকে শুরু হয়েছে; প্রমাণে যে দিকটি
দরকার, সেই ক্রমে লিখি।
\end{quote}
যেহেতু $z \in A$, তাই $z \in A$ অথবা
$z \in B$; সুতরাং~$\cup$-র সংজ্ঞায়
$z \in A \cup B$। একইভাবে $z \in A \cup C$।
অতএব~$\cap$-র সংজ্ঞায়
$z \in (A \cup B) \cap (A \cup C)$।
\begin{quote}
ক্ষেত্রবিচারের প্রথম ক্ষেত্র শেষ। এবার দ্বিতীয়
ক্ষেত্রে, $z \in B \cap C$ হলে, সিদ্ধান্তটি দেখাব।
\end{quote}
ক্ষেত্র 2: ধরো $z \in B \cap C$।
\begin{quote}
আবার দুটি সেটের ছেদ নিয়ে কাজ করছি। $\cap$-র
সংজ্ঞা প্রয়োগ করি:
\end{quote}
যেহেতু $z \in B \cap C$, $\cap$-র সংজ্ঞা
অনুযায়ী $z$ হলো $B$ ও $C$ উভয়েরই
!!a{element}।
\begin{quote}
আবার সিদ্ধান্তটি দেখি। দেখাতে হবে $z$,
$(A \cup B)$ এবং $(A \cup C)$—দুটিতেই আছে।
এবারও সমাধান সরাসরি।
\end{quote}
যেহেতু $z \in B$, তাই $z \in (A \cup B)$।
যেহেতু $z \in C$, তাই $z \in (A
\cup C)$-ও। অতএব $z \in (A \cup B) \cap (A \cup C)$।
\begin{quote}
এখানে আবার $\cup$ ও $\cap$-র সংজ্ঞা প্রয়োগ
করেছি। আগেই সেগুলি বলেছি এবং দেখিয়েছি, দুটি
সেটের একটির সদস্য তাদের সংযোগেরও সদস্য।
তাই এবার ধাপগুলি অতটা স্পষ্ট
করে লেখার দরকার নেই।

ক্ষেত্রবিচারের দ্বিতীয় ক্ষেত্রও শেষ। এখন প্রথম
সিদ্ধান্তটি বলতে পারি।
\end{quote}
অতএব $z \in A \cup (B \cap C)$ হলে
$z \in (A \cup B) \cap (A \cup C)$।
\begin{quote}
এবার শুধু উল্টো দিকটি দেখাতে হবে: $(A \cup B)
\cap (A \cup C)$-র প্রত্যেক !!{element}
$A \cup (B \cap C)$-রও !!a{element}। আগের
মতোই প্রথম সেটের যেকোনো !!{element} নিয়ে
দেখাব, সেটি দ্বিতীয় সেটেও আছে। কী করতে
যাচ্ছি, তা লিখে নিই।
\end{quote}
এবার $z \in (A \cup B) \cap (A \cup C)$ ধরি।
দেখাতে চাই $z \in A \cup (B \cap C)$।
\begin{quote}
এখন পূর্বধারণা $z \in (A \cup B) \cap (A
\cup C)$ থেকে কাজ করছি। প্রমাণের প্রথম অংশের
সেই একই~$z$ আবার ব্যবহার করায় আশা করি
বিভ্রান্তি হবে না। প্রথম অংশ শেষ হলে সেখানকার
সব অনুমান আর কার্যকর থাকে না; তাই এখন $z$
নিয়ে নতুন অনুমান করতে পারি। বিভ্রান্তি লাগলে
পরের অংশে $z$-র বদলে অন্য চলক নাও।

$\cap$-র সংজ্ঞায় $z$ হলো $A \cup B$ ও
$A \cup C$ উভয়ের সদস্য। আবার $\cup$-র
সংজ্ঞা খুললে পাই: $z \in A$ অথবা $z \in B$,
এবং $z \in A$ অথবা $z \in C$। আবার ক্ষেত্রবিচার
মনে হচ্ছে—কিন্তু ``এবং''-টি বিষয়টি জটিল করে।
ভাবতে পারো তিনটি সম্ভাবনা: $z$ হয় $A$, নয় $B$,
নয় $C$-তে আছে। তা ভুল হবে। তাই সাবধানে প্রতিটি
বিযোজন আলাদা করে দেখি।
\end{quote}
$\cap$-র সংজ্ঞায় $z \in A \cup B$ এবং
$z \in A \cup C$। $\cup$-র সংজ্ঞায়
$z \in A$ অথবা $z \in B$। ক্ষেত্রবিচার করি।
\begin{quote}
প্রথম বিযোজনটি নিয়ে ভাবছি বলে দ্বিতীয়টি
($A \cup C$ খুলে পাওয়া) এখনো লিখিনি; আপাতত
তার দরকারও নেই। প্রথম ক্ষেত্র $z \in A$।
কোনো সেটের !!a{element}, সেটটির সঙ্গে অন্য
যেকোনো সেটের সংযোগেরও !!a{element}। তাই
ক্ষেত্র~1 সহজ:
\end{quote}
ক্ষেত্র 1: ধরো $z \in A$। তাহলে
$z \in A \cup (B \cap C)$।
\begin{quote}
এবার দ্বিতীয় ক্ষেত্র $z \in B$। এখানে দ্বিতীয়
$\cup$-র সংজ্ঞা খুলে আবার ক্ষেত্রবিচার করব:
\end{quote}
ক্ষেত্র 2: ধরো $z \in B$। যেহেতু
$z \in A \cup C$, তাই হয় $z \in A$,
নয় $z \in C$। আরও ক্ষেত্রবিচার করি:

ক্ষেত্র 2a: $z \in A$। তাহলে আবার
$z \in A \cup (B \cap C)$।
\begin{quote}
এটি একটু অদ্ভুত: এই ক্ষেত্রে~$z \in B$
পূর্বধারণাটি আসলে লাগেনি। তাতে অসুবিধা নেই।
\end{quote}
ক্ষেত্র 2b: $z \in C$। তাহলে $z \in B$
এবং $z \in C$; তাই $z \in B \cap C$।
ফলে $z \in A \cup (B \cap C)$।
\begin{quote}
উভয় ক্ষেত্রবিচার শেষ; তাই প্রমাণের দ্বিতীয়
অংশও শেষ।
\end{quote}
অতএব $z \in (A \cup B) \cap (A \cup C)$ হলে
$z \in A \cup (B \cap C)$।
\end{proof}

\end{document}
